\chapter{The Bochner Technique}
\label{chap:pet-the-bochner-technique}

Faithful transcription of Petersen, \emph{Riemannian Geometry} (GTM 171, 3rd ed.),
Chapter 9. The Bochner technique starts from Bochner's observation that on a
Riemannian manifold the Bochner--Weitzenb\"ock formula
$\Delta\tfrac12|\nabla u|^2=|\operatorname{Hess}u|^2+\operatorname{Ric}(\nabla u,\nabla u)$
for a harmonic function $u$ (generalizing Bernstein's Euclidean identity
$\Delta\tfrac12|\nabla u|^2=|\operatorname{Hess}u|^2$) carries a curvature term, so that
curvature bounds constrain harmonic functions and, via Hodge's theorem, harmonic
forms and the topology they compute. The chapter develops Hodge theory for the
de Rham complex, the classical Bochner vanishing and P.\ Li--Gallot estimation
theorems for harmonic $1$-forms, the general Lichnerowicz-Laplacian framework
(following an approach due to W.\ A.\ Poor for the Weitzenb\"ock formula on forms)
and its Weitzenb\"ock curvature term, and applications to the curvature tensor,
symmetric $(0,2)$-tensors, and manifolds with nonnegative curvature operator,
using results from earlier sections on Sobolev constants and Moser iteration.

\section{Hodge Theory}

\begin{definition}[de Rham complex and cohomology]
  \label{def:pet-ch9-de-rham-cohomology}
  \lean{PetersenLib.deRhamCohomology}
  On a manifold $M$ the \emph{de Rham complex} is
  \[
    0\to\Omega^0(M)\xrightarrow{d^0}\Omega^1(M)\xrightarrow{d^1}\Omega^2(M)\to\cdots
    \xrightarrow{d^{n-1}}\Omega^n(M)\to0,
  \]
  where $\Omega^k(M)$ is the space of $k$-forms on $M$ and $d^k:\Omega^k(M)\to
  \Omega^{k+1}(M)$ is exterior differentiation. The \emph{de Rham cohomology
  groups} $H^k(M)=\ker(d^k)/\operatorname{im}(d^{k-1})$ compute the real
  cohomology of $M$; $H^0(M)\cong\mathbb{R}$ if $M$ is connected, and $H^n(M)\cong\mathbb{R}$
  if $M$ is orientable and compact.
  \source{petersen:page-0348}
\end{definition}

\begin{remark}[Poincaré duality pairing]
  \label{rem:pet-ch9-poincare-duality}
  \lean{PetersenLib.poincareDualityPairing}
  \uses{def:pet-ch9-de-rham-cohomology}
  When $M$ is compact and oriented the wedge-and-integrate pairing
  $\Omega^k(M)\times\Omega^{n-k}(M)\to\mathbb{R}$, $(\omega_1,\omega_2)\mapsto
  \int_M\omega_1\wedge\omega_2$, induces a nondegenerate pairing on cohomology
  $H^k(M)\times H^{n-k}(M)\to\mathbb{R}$. Hence $H^k(M)$ and $H^{n-k}(M)$ are dual to
  each other and in particular have the same dimension.
  \source{petersen:page-0348,petersen:page-0349}
\end{remark}

\begin{definition}[codifferential and normalized tensor inner product]
  \label{def:pet-ch9-codifferential}
  \lean{PetersenLib.codifferential,PetersenLib.normalizedTensorInnerProduct}
  \uses{def:pet-ch9-de-rham-cohomology}
  When $M$ carries a Riemannian metric $g$, the \emph{codifferential}
  $\delta^k:\Omega^{k+1}(M)\to\Omega^k(M)$ is the adjoint $\nabla^*$ of $d^k$
  characterized by
  \[
    \int_Mg(\delta^k\omega_1,\omega_2)\,\mathrm{vol}=\int_Mg(\omega_1,d^k\omega_2)\,\mathrm{vol}.
  \]
  For tensors of a fixed type one uses the normalized inner product
  $(T_1,T_2)=\tfrac1{\mathrm{vol}\,M}\int_Mg(T_1,T_2)\,\mathrm{vol}$; $\delta$ is also the
  adjoint of $d$ with respect to this normalized pairing.
  \source{petersen:page-0349}
\end{definition}

\begin{definition}[Hodge Laplacian]
  \label{def:pet-ch9-hodge-laplacian}
  \lean{PetersenLib.hodgeLaplacian}
  \uses{def:pet-ch9-codifferential}
  The \emph{Laplacian on forms}, or \emph{Hodge Laplacian}, is
  $\square:\Omega^k(M)\to\Omega^k(M)$, $\square\omega=(d\delta+\delta d)\omega$.
  On functions $\square$ is the negative of the ordinary Laplacian $\Delta$,
  hence the different symbol.
  \source{petersen:page-0349}
\end{definition}

\begin{lemma}[Lemma 9.1.1]
  \label{lem:pet-ch9-harmonic-iff-closed-coclosed}
  \lean{PetersenLib.hodgeLaplacian_eq_zero_iff}
  \uses{def:pet-ch9-hodge-laplacian}
  $\square\omega=0$ if and only if $d\omega=0$ and $\delta\omega=0$. In
  particular $\omega=0$ if $\square\omega=0$ and $\omega=d\theta$ for some
  $\theta$.
  \source{petersen:page-0349}
\end{lemma}
\begin{proof}
  \uses{lem:pet-ch9-harmonic-iff-closed-coclosed}
  Since $d$ and $\delta$ are adjoint,
  \[
    (\square\omega,\omega)=(d\delta\omega,\omega)+(\delta d\omega,\omega)
    =(\delta\omega,\delta\omega)+(d\omega,d\omega).
  \]
  If $\square\omega=0$ this forces $(\delta\omega,\delta\omega)=(d\omega,d\omega)=0$,
  so $\delta\omega=0$ and $d\omega=0$; the converse is immediate from the
  definition of $\square$. If in addition $\omega=d\theta$, then $\delta\omega=0$
  gives $\delta d\theta=0$, i.e.\ $\Delta\theta=0$ (using that $\delta d$ agrees
  with $-\Delta$ on the relevant type), and hence
  $(\omega,\omega)=(d\theta,\omega)=(\theta,\delta\omega)=0$, so $\omega=0$.
\end{proof}

\begin{definition}[Hodge cohomology]
  \label{def:pet-ch9-hodge-cohomology}
  \lean{PetersenLib.hodgeCohomology}
  \uses{def:pet-ch9-hodge-laplacian}
  The \emph{Hodge cohomology} of $(M,g)$ is
  $\mathcal{H}^k(M)=\{\omega\in\Omega^k(M)\mid\square\omega=0\}$, the space of
  \emph{harmonic $k$-forms}.
  \source{petersen:page-0350}
\end{definition}

\begin{theorem}[Hodge, 1935]
  \label{thm:pet-ch9-hodge-theorem}
  \lean{PetersenLib.hodgeCohomology_equiv_deRhamCohomology}
  \uses{def:pet-ch9-hodge-cohomology,lem:pet-ch9-harmonic-iff-closed-coclosed,def:pet-ch9-de-rham-cohomology}
  On a closed oriented Riemannian manifold the natural inclusion
  $\mathcal{H}^k(M)\to H^k(M)$ is an isomorphism.
  \source{petersen:page-0350}
\end{theorem}
\begin{proof}
  \uses{thm:pet-ch9-hodge-theorem,rem:pet-ch9-exercise-9-6-4,rem:pet-ch9-exercise-9-6-5}
  Well-definedness and injectivity of $\mathcal{H}^k(M)\to H^k(M)$ follow from
  \cref{lem:pet-ch9-harmonic-iff-closed-coclosed}. Surjectivity is a standard
  elliptic PDE fact; the spectral theory developed in
  \cref{rem:pet-ch9-exercise-9-6-4,rem:pet-ch9-exercise-9-6-5} establishes part
  of it in a more general context. The key point is that since $\square$ is
  self-adjoint there is an orthogonal decomposition
  \[
    \Omega^k(M)=\operatorname{im}\square\oplus\ker\square
    =\operatorname{im}\square\oplus\mathcal{H}^k(M).
  \]
  Given $\omega\in\Omega^k(M)$, write $\omega=d\delta\theta+\delta d\theta+\tilde\omega$
  with $\square\tilde\omega=0$. If in addition $d\omega=0$, then
  $\square(d\theta)=d\delta(d\theta)=0$; since $d\theta$ is harmonic,
  $\delta(d\theta)=0$ as well by \cref{lem:pet-ch9-harmonic-iff-closed-coclosed}.
  Hence $\omega=d(\delta\theta)+\tilde\omega$, so $\tilde\omega$ represents the
  same cohomology class as $\omega$, proving surjectivity.
\end{proof}

\section{1-Forms}

Hodge theory reduces the computation of the first Betti number
$b_1(M)=\dim\mathcal{H}^1(M)$ to the study of harmonic $1$-forms; the same
program is generalized to other forms and tensors in later sections.

\subsection{The Bochner Formula}

For a harmonic $1$-form $\theta$ on $(M,g)$ let $X$ be the dual vector field,
$\theta(v)=g(X,v)$, and $f=\tfrac12|\theta|^2=\tfrac12|X|^2=\tfrac12\theta(X)$.

\begin{proposition}[Proposition 9.2.1]
  \label{prop:pet-ch9-vector-field-1form-duality}
  \lean{PetersenLib.dualVectorField_symmetric_iff_closed}
  \uses{def:pet-ch9-de-rham-cohomology,def:pet-ch9-codifferential}
  If $X$ and $\theta$ are related by $\theta(v)=g(v,X)$, then
  \begin{enumerate}
    \item $v\mapsto\nabla_vX$ is symmetric if and only if $d\theta=0$, and
    \item $\operatorname{div}X=-\delta\theta$.
  \end{enumerate}
  \source{petersen:page-0350,petersen:page-0351}
\end{proposition}
\begin{proof}
  \uses{prop:pet-ch9-vector-field-1form-duality}
  Recall the identity $d\theta(V,W)+(\mathcal{L}_Xg)(V,W)=2g(\nabla_VX,W)$. Since
  $\mathcal{L}_Xg$ is symmetric in $V,W$ and $d\theta$ is skew-symmetric, comparing
  symmetric and skew-symmetric parts of $2g(\nabla_VX,W)$ shows $v\mapsto\nabla_vX$
  is symmetric exactly when $d\theta=0$. The second identity,
  $\operatorname{div}X=-\delta\theta$, is the $1$-form instance of the general duality
  between divergence and codifferential (Petersen, Proposition 2.2.7).
\end{proof}

Consequently, when $\theta$ is harmonic then $\operatorname{div}X=0$ and $\nabla X$ is a
symmetric $(1,1)$-tensor; the Bochner formula for closed $1$-forms is stated
below through the dual vector field.

\begin{proposition}[Proposition 9.2.2]
  \label{prop:pet-ch9-bochner-formula-1forms}
  \lean{PetersenLib.bochnerFormula_oneForm}
  \uses{prop:pet-ch9-vector-field-1form-duality}
  Let $X$ be a vector field with $\nabla X$ symmetric (equivalently the dual
  $1$-form is closed), $f=\tfrac12|X|^2$, and suppose $X=\nabla u$ near $p$. Then
  \begin{enumerate}
    \item $\nabla f=\nabla_XX$.
    \item $\operatorname{Hess}f(V,V)=\operatorname{Hess}^2u(V,V)+(\nabla_X\operatorname{Hess}u)(V,V)+R(V,X,X,V)
      =|\nabla_VX|^2+g(\nabla^2_XX,V)+R(V,X,X,V)$.
    \item $\Delta f=|\operatorname{Hess}u|^2+D_X\Delta u+\operatorname{Ric}(X,X)
      =|\nabla X|^2+D_X\operatorname{div}X+\operatorname{Ric}(X,X)$.
  \end{enumerate}
  \source{petersen:page-0351}
\end{proposition}
\begin{proof}
  \uses{prop:pet-ch9-bochner-formula-1forms}
  For (1), $g(\nabla f,V)=D_V\tfrac12|X|^2=g(\nabla_VX,X)=g(\nabla_XX,V)$, using
  symmetry of $\nabla X$ in the last step. Part (2) is theorem 3.2.2 of the
  source applied to $u$. For (3), take traces in (2): as in the proof of
  Petersen's proposition 8.2.1, the trace gives the first and third terms
  directly; the second comes from commuting the trace with a covariant
  derivative. At a point $p$, either $X|_p=0$ or one may choose an orthonormal
  frame $E_i$ parallel along $X$ at $p$; in either case
  \[
    \sum_ig(\nabla^2_{E_i,X}X,E_i)=\sum_i(\nabla_X\operatorname{Hess}u)(E_i,E_i)
    =D_X\sum_i\operatorname{Hess}u(E_i,E_i)=D_X\Delta u,
  \]
  which is the claimed second term.
\end{proof}

\subsection{The Vanishing Theorem}

\begin{theorem}[Theorem 9.2.3 (Bochner, 1948)]
  \label{thm:pet-ch9-bochner-vanishing-1forms}
  \lean{PetersenLib.bochner_harmonicOneForm_parallel}
  \uses{prop:pet-ch9-bochner-formula-1forms}
  If $(M,g)$ is compact and $\operatorname{Ric}\ge0$, then every harmonic $1$-form is
  parallel.
  \source{petersen:page-0352}
\end{theorem}
\begin{proof}
  \uses{thm:pet-ch9-bochner-vanishing-1forms}
  Let $\omega$ be a harmonic $1$-form and $X$ its dual vector field. Since
  $\omega$ is closed and coclosed, $\nabla X$ is symmetric and
  $\operatorname{div}X=\Delta u=0$ (in the notation of
  \cref{prop:pet-ch9-bochner-formula-1forms}), so part (3) of that proposition
  gives
  \[
    \Delta\Big(\tfrac12|X|^2\Big)=|\nabla X|^2+\operatorname{Ric}(X,X)\ge0.
  \]
  On a compact manifold the maximum principle applied to the subharmonic
  function $\tfrac12|X|^2$ forces it to be constant, and then equality
  throughout forces $|\nabla X|=0$, i.e.\ $X$ (equivalently $\omega$) is
  parallel.
\end{proof}

\begin{corollary}[Corollary 9.2.4]
  \label{cor:pet-ch9-harmonic-1forms-vanish-positive-ricci}
  \lean{PetersenLib.harmonicOneForms_vanish_of_positiveRicci}
  \uses{thm:pet-ch9-bochner-vanishing-1forms}
  If $(M,g)$ is as in \cref{thm:pet-ch9-bochner-vanishing-1forms} and moreover
  $\operatorname{Ric}$ is positive at one point, then every harmonic $1$-form vanishes
  identically.
  \source{petersen:page-0352}
\end{corollary}
\begin{proof}
  \uses{cor:pet-ch9-harmonic-1forms-vanish-positive-ricci}
  By \cref{thm:pet-ch9-bochner-vanishing-1forms}, $\operatorname{Ric}(X,X)\equiv0$ for the
  dual vector field $X$ of a harmonic $1$-form. If $\operatorname{Ric}$ is positive
  definite at $p$, this forces $X|_p=0$; since $X$ is parallel, $X\equiv0$.
\end{proof}

\begin{corollary}[Corollary 9.2.5]
  \label{cor:pet-ch9-betti1-bound-flat-torus}
  \lean{PetersenLib.firstBettiNumber_le_dim_iff_flatTorus}
  \uses{thm:pet-ch9-bochner-vanishing-1forms,def:pet-ch9-hodge-cohomology}
  If $(M,g)$ is compact with $\operatorname{Ric}\ge0$, then $b_1(M)\le n=\dim M$, with
  equality if and only if $(M,g)$ is a flat torus.
  \source{petersen:page-0352}
\end{corollary}
\begin{proof}
  \uses{cor:pet-ch9-betti1-bound-flat-torus}
  By Hodge theory $b_1(M)=\dim\mathcal{H}^1(M)$, and by
  \cref{thm:pet-ch9-bochner-vanishing-1forms} every harmonic $1$-form is
  parallel, so the evaluation map $\mathcal{H}^1(M)\to T_p^*M$ at any $p$ is
  injective; hence $\dim\mathcal{H}^1(M)\le n$.

  If equality holds, there are $n$ linearly independent parallel vector fields
  $E_1,\dots,E_n$, forcing $(M,g)$ to be flat. The universal cover is then
  $\mathbb{R}^n$ with $\pi_1(M)$ acting by isometries; pulling the $E_i$ back to
  $\mathbb{R}^n$ gives constant vector fields, which may be taken to be the standard
  coordinate fields $\partial_i$. Invariance of the $E_i$ under $\pi_1(M)$ means
  that for every $F\in\pi_1(M)$, $DF(\partial_i|_p)=\partial_i|_{F(p)}$ for all $i$;
  only translations fix every coordinate field, so $\pi_1(M)$ consists entirely
  of translations. Thus $\pi_1(M)$ is finitely generated, abelian, and torsion
  free, so $\Gamma:=\pi_1(M)\cong\mathbb{Z}^k$ for some $k$. If $k<n$, $\mathbb{Z}^k$
  generates a proper subspace of translations and cannot act cocompactly on
  $\mathbb{R}^n$; if $k>n$, $\mathbb{Z}^k$ cannot act discretely on $\mathbb{R}^n$. Hence $k=n$,
  $\Gamma\cong\mathbb{Z}^n$ generates $\mathbb{R}^n$, and $M=\mathbb{R}^n/\Gamma$ is a flat torus.
  Conversely a flat torus has $\mathbb{R}^n$ worth of parallel $1$-forms, so
  $b_1=n$.
\end{proof}

\subsection{The Estimation Theorem}

The goal is to generalize \cref{thm:pet-ch9-bochner-vanishing-1forms} to
manifolds with a negative lower Ricci bound, following P.\ Li's technique as
sharpened by Gallot, whose contribution included a suitable bound on Sobolev
constants (Petersen, theorem 7.1.13).

\begin{lemma}[Lemma 9.2.6 (P.\ Li)]
  \label{lem:pet-ch9-li-dimension-estimate}
  \lean{PetersenLib.li_dimension_estimate}
  Let $(M,g)$ be compact and $E\to M$ a vector bundle of fibre dimension $m$
  with a smoothly varying fibre inner product. For $s\in\Gamma(E)$ set
  $\|s\|_\infty=\max_{x\in M}|s(x)|$ and $\|s\|_p=\big(\tfrac1{\mathrm{vol}\,M}
  \int_M|s|^p\,\mathrm{vol}\big)^{1/p}$. For a finite-dimensional $V\subset\Gamma(E)$
  put $C(V)=\max_{s\in V\setminus\{0\}}\|s\|_\infty/\|s\|_2$. Then
  $\dim V\le m\cdot C(V)$.
  \source{petersen:page-0353,petersen:page-0354}
\end{lemma}
\begin{proof}
  \uses{lem:pet-ch9-li-dimension-estimate}
  $V$ carries the inner product $(s_1,s_2)=\tfrac1{\mathrm{vol}\,M}\int_M
  \langle s_1,s_2\rangle\,\mathrm{vol}$ with $(s,s)=\|s\|_2^2$. Choose an orthonormal
  basis $e_1,\dots,e_\ell$ of $V$ for this inner product and set
  $f(x)=\sum_{i=1}^\ell|e_i(x)|^2$, a function independent of the choice of
  orthonormal basis. Averaging term by term,
  \[
    \frac1{\mathrm{vol}\,M}\int_Mf\,\mathrm{vol}=\ell=\dim V.
  \]
  Let $x_0$ maximize $f$. The evaluation map $V\to E_{x_0}$ has rank $k\le m$
  (its target has dimension $m$), so the orthonormal basis may be chosen so
  that the last $\ell-k$ elements span its kernel, i.e.\ vanish at $x_0$; thus
  $f(x_0)=\sum_{i=1}^k|e_i(x_0)|^2$. Since $\|e_i\|_2=1$ for each $i$,
  $|e_i(x_0)|\le\|e_i\|_\infty\le C(V)\|e_i\|_2=C(V)$, and summing the $k$
  nonvanishing terms gives $f(x_0)\le k\cdot C(V)\le m\cdot C(V)$. As the
  average of $f$ is at most its maximum,
  \[
    \dim V=\ell\le f(x_0)\le m\cdot C(V).\qedhere
  \]
\end{proof}

\begin{theorem}[Theorem 9.2.7 (Moser iteration)]
  \label{thm:pet-ch9-moser-iteration}
  \lean{PetersenLib.moserIteration}
  Let $(M,g)$ be compact and suppose $\|u\|_{2\nu}\le S\|\nabla u\|_2+\|u\|_2$
  for all smooth $u$ on $M$, with $\nu>1$. If $f:M\to[0,\infty)$ is continuous,
  smooth on $\{f>0\}$, and $\Delta f\ge-\lambda f$, then
  \[
    \|f\|_\infty\le\exp\Big(\frac{S\sqrt{\lambda\nu}}{\sqrt{\nu-1}}\Big)\|f\|_2.
  \]
  \source{petersen:page-0354,petersen:page-0355}
\end{theorem}
\begin{proof}
  \uses{thm:pet-ch9-moser-iteration}
  Since $f$ is minimized (namely $=0$) on $\{f=0\}$, all its derivatives vanish
  there, and $\Delta f\ge0$ there in both the barrier and distributional
  sense; hence the differential inequality $\Delta f\ge-\lambda f$ holds
  distributionally on all of $M$. Green's formula gives, for $q\ge1$,
  \[
    (f^{2q-1},\Delta f)=-(df^{2q-1},df)=-(2q-1)(f^{2q-2}df,df),
  \]
  so
  \[
    \|df^q\|_2^2=q^2(f^{2q-2}df,df)=-\frac{q^2}{2q-1}(f^{2q-1},\Delta f)
    \le\frac{q^2\lambda}{2q-1}(f^{2q-1},f)=\frac{q^2\lambda}{2q-1}\|f^q\|_2^2.
  \]
  Applying the Sobolev inequality to $f^q$,
  \[
    \|f^q\|_{2\nu}\le S\|df^q\|_2+\|f^q\|_2
    \le\Big(Sq\Big(\frac\lambda{2q-1}\Big)^{1/2}+1\Big)\|f^q\|_2,
  \]
  equivalently
  \[
    \|f\|_{2q\nu}\le\Big(Sq\Big(\frac\lambda{2q-1}\Big)^{1/2}+1\Big)^{1/q}\|f\|_{2q}.
  \]
  Setting $q=\nu^k$,
  \[
    \|f\|_{2\nu^{k+1}}\le\Big(S\nu^k\Big(\frac\lambda{2\nu^k-1}\Big)^{1/2}+1\Big)^{\nu^{-k}}\|f\|_{2\nu^k}.
  \]
  Iterating from $k=0$ to $\infty$,
  \[
    \|f\|_\infty\le\Big(\prod_{k=0}^\infty\Big(S\nu^k\Big(\frac\lambda{2\nu^k-1}\Big)^{1/2}+1\Big)^{\nu^{-k}}\Big)\|f\|_2.
  \]
  Taking logarithms and using $\log(1+x)\le x$,
  \[
    \sum_{k=0}^\infty\nu^{-k}\log\Big(S\nu^k\Big(\frac\lambda{2\nu^k-1}\Big)^{1/2}+1\Big)
    \le S\sqrt\lambda\sum_{k=0}^\infty\Big(\frac1{2\nu^k-1}\Big)^{1/2}
    \le S\sqrt\lambda\sum_{k=0}^\infty\Big(\frac1{\nu^k}\Big)^{1/2}
    =\frac{S\sqrt{\lambda\nu}}{\sqrt{\nu-1}},
  \]
  which exponentiates to the claimed bound.
\end{proof}

\begin{lemma}[Kato's inequality]
  \label{lem:pet-ch9-kato-inequality}
  \lean{PetersenLib.katoInequality}
  For a tensor field $T$ with $f=|T|$, smooth away from $\{T=0\}$, one has
  $df\le|\nabla T|$ pointwise wherever $f$ is smooth.
  \source{petersen:page-0356}
\end{lemma}
\begin{proof}
  \uses{lem:pet-ch9-kato-inequality}
  Where $f>0$, $2f\,df=df^2=2g(\nabla T,T)\le2|\nabla T|f$ by Cauchy--Schwarz,
  and dividing by $2f$ gives $df\le|\nabla T|$.
\end{proof}

\begin{theorem}[Theorem 9.2.8 (Gromov, 1980 and Gallot, 1981)]
  \label{thm:pet-ch9-gromov-gallot-betti1-estimate}
  \lean{PetersenLib.gromovGallot_firstBettiNumber_estimate}
  \uses{lem:pet-ch9-li-dimension-estimate,thm:pet-ch9-moser-iteration,lem:pet-ch9-kato-inequality,def:pet-ch9-hodge-cohomology,prop:pet-ch9-bochner-formula-1forms}
  If $M$ is a compact Riemannian $n$-manifold with $\operatorname{Ric}\ge(n-1)k$ and
  $\operatorname{diam}(M)\le D$, there is a function $C(n,kD^2)$ with
  $b_1(M)\le C(n,kD^2)$ and $\lim_{\varepsilon\to0}C(n,\varepsilon)=n$. In
  particular there is $\varepsilon(n)>0$ so that $kD^2\ge-\varepsilon(n)$
  implies $b_1(M)\le n$.
  \source{petersen:page-0355,petersen:page-0356}
\end{theorem}
\begin{proof}
  \uses{thm:pet-ch9-gromov-gallot-betti1-estimate}
  By \cref{lem:pet-ch9-li-dimension-estimate} applied to
  $V=\mathcal{H}^1(M)$ (fibre dimension $m=n$),
  $\dim\mathcal{H}^1\le n\cdot C(\mathcal{H}^1)$, so it suffices to bound the
  ratios $\|\omega\|_\infty/\|\omega\|_2$ for harmonic $1$-forms $\omega$. Let
  $f=|\omega|$ and $X$ the vector field dual to $\omega$. By
  \cref{lem:pet-ch9-kato-inequality}, $df\le|\nabla\omega|$. By the Bochner
  formula (\cref{prop:pet-ch9-bochner-formula-1forms} applied with
  $\operatorname{div}X=0$),
  \[
    |df|^2+f\Delta f=\tfrac12\Delta f^2=|\nabla\omega|^2+\operatorname{Ric}(X,X)
    \ge|\nabla\omega|^2+(n-1)kf^2,
  \]
  and combining with Kato's inequality gives $\Delta f\ge(n-1)kf$. Theorem
  \ref{thm:pet-ch9-moser-iteration} (applied with $\lambda=-(n-1)k$ when
  $k\le0$) then yields
  \[
    \|f\|_\infty\le\exp\Big(\frac{S\sqrt{-(n-1)k\nu}}{\sqrt{\nu-1}}\Big)\|f\|_2,
  \]
  where $S=D\cdot C(n,kD^2)$ is the Sobolev constant of Petersen's theorem
  7.1.13 and proposition 7.1.17. Since $\|\omega\|_\infty/\|\omega\|_2=
  \|f\|_\infty/\|f\|_2$,
  \[
    \dim\mathcal{H}^1\le n\cdot C(\mathcal{H}^1)
    \le n\cdot\exp\Big(\frac{S\sqrt{-(n-1)k\nu}}{\sqrt{\nu-1}}\Big),
  \]
  which is a function of $n$ and $kD^2$ alone, proving the theorem; as
  $kD^2\to0$ the bound tends to $n$.
\end{proof}

\section{Lichnerowicz Laplacians}

A natural class of Laplacians is introduced, generalizing the Bochner
technique used above for $1$-forms; §9.4 exhibits several natural examples,
including the Hodge Laplacian.

\subsection{The Connection Laplacian}

Fix a tensor bundle $E\to M$ of fibre dimension $m$ (e.g.\ $p$-forms,
symmetric tensors, or curvature tensors).

\begin{proposition}[Proposition 9.3.1]
  \label{prop:pet-ch9-vanishing-connection-laplacian}
  \lean{PetersenLib.vanishing_connectionLaplacian}
  Let $(M,g)$ be Riemannian and $T\in\Gamma(E)$ with $g(\nabla^*\nabla T,T)\le0$.
  If $|T|$ attains a maximum, then $T$ is parallel.
  \source{petersen:page-0357}
\end{proposition}
\begin{proof}
  \uses{prop:pet-ch9-vanishing-connection-laplacian}
  One has $\Delta\tfrac12|T|^2=|\nabla T|^2-g(\nabla^*\nabla T,T)\ge0$. If $|T|$
  has a maximum, the maximum principle applied to $|T|^2$ shows it is
  constant, forcing $|\nabla T|\equiv0$, i.e.\ $T$ is parallel.
\end{proof}

\begin{theorem}[Theorem 9.3.2 (Gallot, 1981)]
  \label{thm:pet-ch9-gallot-connection-laplacian-estimate}
  \lean{PetersenLib.gallot_connectionLaplacian_estimate}
  \uses{lem:pet-ch9-li-dimension-estimate,thm:pet-ch9-moser-iteration,lem:pet-ch9-kato-inequality}
  Assume $(M,g)$ is compact and satisfies the hypothesis of
  \cref{thm:pet-ch9-moser-iteration}. Let $V\subset\Gamma(E)$ be finite
  dimensional. If $g(\nabla^*\nabla T,T)\le\lambda|T|^2$ for all $T\in V$, then
  \[
    \dim V\le m\cdot\exp\Big(\frac{S\sqrt\lambda}{\sqrt{\nu-1}}\Big).
  \]
  \source{petersen:page-0357}
\end{theorem}
\begin{proof}
  \uses{thm:pet-ch9-gallot-connection-laplacian-estimate}
  As in \cref{thm:pet-ch9-gromov-gallot-betti1-estimate}, set $f=|T|$ for
  $T\in V$. Instead of the $1$-form Bochner formula one uses
  \[
    |df|^2+f\Delta f=\tfrac12\Delta f^2=|\nabla T|^2-g(\nabla^*\nabla T,T)\ge0,
  \]
  together with $g(\nabla^*\nabla T,T)\le\lambda|T|^2$, to get $f\Delta f\ge
  -\lambda f^2$; by \cref{lem:pet-ch9-kato-inequality}, $\Delta f\ge-\lambda f$.
  Theorem \ref{thm:pet-ch9-moser-iteration} then gives $\|f\|_\infty\le
  \exp(S\sqrt\lambda/\sqrt{\nu-1})\|f\|_2$ for every $T\in V$, so
  $C(V)\le\exp(S\sqrt\lambda/\sqrt{\nu-1})$, and
  \cref{lem:pet-ch9-li-dimension-estimate} finishes the proof.
\end{proof}

\subsection{The Weitzenböck Curvature}

\begin{definition}[Weitzenböck curvature operator]
  \label{def:pet-ch9-weitzenbock-curvature-operator}
  \lean{PetersenLib.weitzenbockCurvatureOperator}
  For a $(0,k)$-tensor $T$ the \emph{Weitzenböck curvature} of $T$ is
  \[
    \operatorname{Ric}(T)(X_1,\dots,X_k)=\sum_j\big(R(e_j,X_i)T\big)(X_1,\dots,e_j,\dots,X_k)
  \]
  (sum over $i$ implicit in the slot notation, and over an orthonormal frame
  $e_j$). For a $1$-form $\omega$ dual to a vector field $X$, $\operatorname{Ric}(\omega)(X')
  =\sum_j(R(e_j,X')\omega)(e_j)=-\omega\big(\sum_jR(e_j,X')e_j\big)=\omega(\operatorname{Ric}(X'))$,
  recovering the ordinary Ricci tensor; the operator is often written $W$, but
  $\operatorname{Ric}$ is used here to stress this compatibility.
  \source{petersen:page-0357,petersen:page-0358}
\end{definition}

\begin{definition}[Lichnerowicz Laplacian]
  \label{def:pet-ch9-lichnerowicz-laplacian}
  \lean{PetersenLib.lichnerowiczLaplacian}
  \uses{def:pet-ch9-weitzenbock-curvature-operator}
  For a constant $c>0$, the \emph{Lichnerowicz Laplacian} on sections of $E$ is
  $\Delta_LT=\nabla^*\nabla T+c\operatorname{Ric}(T)$. The Hodge Laplacian on forms is of
  this type with $c=1$; symmetric $(0,2)$-tensors and the curvature tensor use
  $c=\tfrac12$.
  \source{petersen:page-0358}
\end{definition}

The Bochner technique applies to $T$ with $\Delta_LT=0$: the maximum principle
(\cref{prop:pet-ch9-vanishing-connection-laplacian}) forces $T$ parallel once
$g(\nabla^*\nabla T,T)\le0$, which by $\Delta_LT=0$ is equivalent to
$g(\operatorname{Ric}(T),T)\ge0$. The remaining work is to reformulate
$g(\operatorname{Ric}(T),T)$ so that nonnegativity of the curvature operator makes this
manifest, and to identify natural Lichnerowicz Laplacians.

\subsection{Simplification of $\operatorname{Ric}(T)$}

\begin{convention}[metric on $\mathfrak{so}(T_pM)$]
  \label{conv:pet-ch9-so-metric-convention}
  \lean{PetersenLib.soMetricConvention}
  \uses{def:pet-ch9-weitzenbock-curvature-operator}
  For each $X,Y\in T_pM$ the map $R_{X,Y}:T_pM\to T_pM$ is skew-symmetric,
  hence decomposes in an orthonormal basis $\Xi_\sigma$ of
  $\mathfrak{so}(T_pM)$. Since an orthonormal $v,w$ give a unit bivector
  $v\wedge w\in\Lambda^2T_pM$ while the corresponding rotation by $\pi/2$ has
  Euclidean operator norm $\sqrt2$, $\mathfrak{so}(T_pM)$ is endowed with the
  metric transported from $\Lambda^2T_pM$ under $v\wedge w\leftrightarrow$
  (rotation by $\pi/2$ in $\operatorname{span}\{v,w\}$). With the curvature operator
  $\mathfrak{R}$ on $\Lambda^2T_pM\cong\mathfrak{so}(T_pM)$ this gives
  \[
    R_{X,Y}=\sum_\sigma g(R_{X,Y},\Xi_\sigma)\Xi_\sigma
    =\sum_\sigma g(\mathfrak{R}(X\wedge Y),\Xi_\sigma)\Xi_\sigma
    =\sum_\sigma g(\mathfrak{R}(\Xi_\sigma),X\wedge Y)\Xi_\sigma
    =-\sum_\sigma g(R(\Xi_\sigma)X,Y)\Xi_\sigma.
  \]
  \source{petersen:page-0359}
\end{convention}

\begin{lemma}[Lemma 9.3.3]
  \label{lem:pet-ch9-ric-curvature-operator-formula}
  \lean{PetersenLib.weitzenbock_eq_curvatureOperator_sum}
  \uses{conv:pet-ch9-so-metric-convention,def:pet-ch9-weitzenbock-curvature-operator,def:pet-ch9-lichnerowicz-laplacian}
  For any $(0,k)$-tensor $T$,
  \[
    \operatorname{Ric}(T)=-\sum_\sigma R(\Xi_\sigma)(\Xi_\sigma T),\qquad
    \Delta_LT=\nabla^*\nabla T-c\sum_\sigma R(\Xi_\sigma)(\Xi_\sigma T).
  \]
  Moreover $\operatorname{Ric}$ is self-adjoint.
  \source{petersen:page-0359,petersen:page-0360}
\end{lemma}
\begin{proof}
  \uses{lem:pet-ch9-ric-curvature-operator-formula}
  Using \cref{conv:pet-ch9-so-metric-convention} to expand $R(e_j,X_i)$,
  \begin{align*}
    \operatorname{Ric}(T)(X_1,\dots,X_k)
    &=\sum_j(R(e_j,X_i)T)(X_1,\dots,e_j,\dots,X_k)\\
    &=-\sum_{\sigma,j}g(R(\Xi_\sigma)e_j,X_i)(\Xi_\sigma T)(X_1,\dots,e_j,\dots,X_k)\\
    &=-\sum_\sigma(\Xi_\sigma T)\big(X_1,\dots,g(R(\Xi_\sigma)e_j,X_i)e_j,\dots,X_k\big)\\
    &=\sum_\sigma(\Xi_\sigma T)(X_1,\dots,R(\Xi_\sigma)X_i,\dots,X_k)\\
    &=-\sum_\sigma\big(R(\Xi_\sigma)(\Xi_\sigma T)\big)(X_1,\dots,X_i,\dots,X_k),
  \end{align*}
  where the fourth equality resums the $e_j$-expansion of $R(\Xi_\sigma)X_i$
  back into the slot, and the fifth uses the definition of $R(\Xi_\sigma)$
  acting on tensors by derivation through the slot $X_i$. Substituting into
  $\Delta_LT=\nabla^*\nabla T+c\operatorname{Ric}(T)$ gives the stated formula for
  $\Delta_L$.

  For self-adjointness, choose the orthonormal basis $\Xi_\sigma$ to diagonalize
  $\mathfrak{R}$, $\mathfrak{R}(\Xi_\sigma)=\lambda_\sigma\Xi_\sigma$. Then
  \[
    g(\operatorname{Ric}(T),S)=-\sum_\sigma g(R(\Xi_\sigma)(\Xi_\sigma T),S)
    =-\sum_\sigma\lambda_\sigma g(\Xi_\sigma(\Xi_\sigma T),S)
    =\sum_\sigma\lambda_\sigma g(\Xi_\sigma T,\Xi_\sigma S),
  \]
  using that $\Xi_\sigma$ acts skew-adjointly on tensors, and this final
  expression is symmetric in $T,S$.
\end{proof}

\begin{corollary}[Corollary 9.3.4]
  \label{cor:pet-ch9-ric-nonneg-curvature-operator}
  \lean{PetersenLib.weitzenbock_nonneg_of_curvatureOperator_nonneg}
  \uses{lem:pet-ch9-ric-curvature-operator-formula}
  If $\mathfrak{R}\ge0$, then $g(\operatorname{Ric}(T),T)\ge0$ for every tensor $T$. More
  generally, if $\mathfrak{R}\ge k$ with $k<0$, then $g(\operatorname{Ric}(T),T)\ge kC|T|^2$
  for a constant $C>0$ depending only on the tensor type.
  \source{petersen:page-0360}
\end{corollary}
\begin{proof}
  \uses{cor:pet-ch9-ric-nonneg-curvature-operator}
  With $\mathfrak{R}(\Xi_\sigma)=\lambda_\sigma\Xi_\sigma$ as in the proof of
  \cref{lem:pet-ch9-ric-curvature-operator-formula},
  \[
    g(\operatorname{Ric}(T),T)=\sum_\sigma\lambda_\sigma|\Xi_\sigma T|^2\ge k\sum_\sigma|\Xi_\sigma T|^2,
  \]
  which is $\ge0$ when $k=0$. There is a constant $C>0$, depending only on the
  tensor type and $\dim M$, with $\sum_\sigma|\Xi_\sigma T|^2\le C|T|^2$; when
  $k<0$ this yields $g(\operatorname{Ric}(T),T)\ge kC|T|^2$.
\end{proof}

\begin{theorem}[Theorem 9.3.5]
  \label{thm:pet-ch9-lichnerowicz-dimension-estimate}
  \lean{PetersenLib.lichnerowiczLaplacian_kernel_dimension_estimate}
  \uses{cor:pet-ch9-ric-nonneg-curvature-operator,thm:pet-ch9-gallot-connection-laplacian-estimate,prop:pet-ch9-vanishing-connection-laplacian}
  If $\mathfrak{R}\ge k$ and $\operatorname{diam}(M)\le D$, the dimension of
  $V=\{T\in\Gamma(E)\mid\Delta_LT=\nabla^*\nabla T+c\operatorname{Ric}(T)=0\}$ is bounded by
  \[
    m\cdot\exp\Big(D\cdot C(n,kD^2)\frac{\sqrt{kcC\nu}}{\sqrt{\nu-1}}\Big),
  \]
  and when $k=0$ every $T\in V$ is parallel.
  \source{petersen:page-0360}
\end{theorem}

\section{The Bochner Technique in General}

Several natural Lichnerowicz Laplacians on Riemannian manifolds are exhibited.

\subsection{Forms}

\begin{theorem}[Theorem 9.4.1 (Weitzenböck, 1923)]
  \label{thm:pet-ch9-weitzenbock-hodge-laplacian}
  \lean{PetersenLib.hodgeLaplacian_eq_lichnerowiczLaplacian}
  \uses{def:pet-ch9-hodge-laplacian,def:pet-ch9-weitzenbock-curvature-operator,def:pet-ch9-lichnerowicz-laplacian}
  The Hodge Laplacian is the Lichnerowicz Laplacian with $c=1$:
  \[
    \square\omega=(d\delta+\delta d)\omega=\nabla^*\nabla\omega+\operatorname{Ric}(\omega).
  \]
  \source{petersen:page-0361}
\end{theorem}
\begin{proof}
  \uses{thm:pet-ch9-weitzenbock-hodge-laplacian}
  Following W.\ A.\ Poor's argument, write, for a $k$-form $\omega$ and an
  orthonormal frame $E_j$,
  \begin{align*}
    \delta\omega(X_2,\dots,X_k)&=-\sum_i(\nabla_{E_i}\omega)(E_i,X_2,\dots,X_k),\\
    d\omega(X_0,\dots,X_k)&=\sum_i(-1)^i(\nabla_{X_i}\omega)(X_0,\dots,\hat X_i,\dots,X_k).
  \end{align*}
  Substituting the first into $d$ and the second into $\delta$ (and
  differentiating again) gives, after re-indexing the alternating sums,
  \begin{align*}
    d\delta\omega(X_1,\dots,X_k)
    &=-\sum_i(-1)^{i+1}(\nabla^2_{X_i,E_j}\omega)(E_j,X_1,\dots,\hat X_i,\dots,X_k)\\
    &=-\sum_{i,j}(\nabla^2_{X_i,E_j}\omega)(X_1,\dots,E_j,\dots,X_k),\\[4pt]
    \delta d\omega(X_1,\dots,X_k)
    &=-\sum_j(\nabla^2_{E_j,E_j}\omega)(X_1,\dots,X_k)
      -\sum_i(-1)^i(\nabla^2_{E_j,X_i}\omega)(E_j,X_1,\dots,\hat X_i,\dots,X_k)\\
    &=(\nabla^*\nabla\omega)(X_1,\dots,X_k)
      +\sum_{i,j}(\nabla^2_{E_j,X_i}\omega)(X_1,\dots,E_j,\dots,X_k).
  \end{align*}
  Adding the two displays and using the second Bianchi/curvature commutation
  identity $\nabla^2_{E_j,X_i}\omega-\nabla^2_{X_i,E_j}\omega=R(E_j,X_i)\omega$
  slot by slot,
  \[
    \square\omega=\nabla^*\nabla\omega+\sum_{i,j}\big(R(E_j,X_i)\omega\big)(X_1,\dots,E_j,\dots,X_k)
    =\nabla^*\nabla\omega+\operatorname{Ric}(\omega).\qedhere
  \]
\end{proof}

\subsection{The Curvature Tensor}

\begin{theorem}[Theorem 9.4.2]
  \label{thm:pet-ch9-curvature-tensor-lichnerowicz}
  \lean{PetersenLib.curvatureTensor_lichnerowiczIdentity}
  \uses{def:pet-ch9-lichnerowicz-laplacian,def:pet-ch9-weitzenbock-curvature-operator}
  The curvature tensor $R$ of a Riemannian manifold satisfies
  \begin{align*}
    (\nabla^*\nabla R)(X,Y,Z,W)+\tfrac12\operatorname{Ric}(R)(X,Y,Z,W)
    &=\tfrac12(\nabla_X\nabla^*R)(Y,Z,W)-\tfrac12(\nabla_Y\nabla^*R)(X,Z,W)\\
    &\quad+\tfrac12(\nabla_Z\nabla^*R)(W,X,Y)-\tfrac12(\nabla_W\nabla^*R)(Z,X,Y).
  \end{align*}
  In particular $\Delta_LR=\nabla^*\nabla R+\tfrac12\operatorname{Ric}(R)$ vanishes whenever
  $\nabla^*R=0$, i.e.\ $R$ is a Lichnerowicz-harmonic tensor ($c=\tfrac12$)
  precisely when it is divergence free.
  \source{petersen:page-0362,petersen:page-0363}
\end{theorem}
\begin{proof}
  \uses{thm:pet-ch9-curvature-tensor-lichnerowicz}
  The second Bianchi identity is the key ingredient. Fix $p\in M$, let
  $X,Y,Z,W$ be vector fields with $\nabla X=\nabla Y=\nabla Z=\nabla W=0$ at
  $p$, and let $E_i$ be a normal frame at $p$. Then, evaluating at $p$,
  \begin{align*}
    (\nabla^*\nabla R)(X,Y,Z,W)
    &=-\sum_i(\nabla^2_{E_i,E_i}R)(X,Y,Z,W)\\
    &=\sum_i(\nabla^2_{E_i,X}R)(Y,E_i,Z,W)+(\nabla^2_{E_i,Y}R)(E_i,X,Z,W)\\
    &=\sum_i(\nabla^2_{X,E_i}R)(Y,E_i,Z,W)+(\nabla^2_{Y,E_i}R)(E_i,X,Z,W)\\
    &\quad+\sum_i\big(R(E_i,X)(R)\big)(Y,E_i,Z,W)+\big(R(E_i,Y)(R)\big)(E_i,X,Z,W)\\
    &=(\nabla_X\nabla^*R)(Y,Z,W)-(\nabla_Y\nabla^*R)(X,Z,W)\\
    &\quad-\sum_i\big(R(E_i,X)(R)\big)(E_i,Y,Z,W)+\big(R(E_i,Y)(R)\big)(X,E_i,Z,W),
  \end{align*}
  where the second line uses the second Bianchi identity to rewrite the
  divergence of $R$ against $E_i$ in terms of derivatives of the entries $X,Y$
  (which are parallel at $p$, so covariant derivatives of $R$ against them may
  be pulled past the sum over $i$), the third line inserts the curvature
  commutation $\nabla^2_{E_i,X}-\nabla^2_{X,E_i}=R(E_i,X)$ acting on $R$ (using
  $\nabla X=0$ at $p$), and the fourth resums $\sum_i(\nabla^2_{X,E_i}R)(Y,E_i,Z,W)
  =(\nabla_X\nabla^*R)(Y,Z,W)$ by definition of $\nabla^*R$. The last two terms
  are half of $-\operatorname{Ric}(R)(X,Y,Z,W)$, i.e.
  \[
    -\sum_i\big(R(E_i,X)(R)\big)(E_i,Y,Z,W)+\big(R(E_i,Y)(R)\big)(X,E_i,Z,W)
    =-\tfrac12\operatorname{Ric}(R)(X,Y,Z,W)+(\ast),
  \]
  where $(\ast)$ collects the analogous pair-symmetric terms with the roles of
  $(X,Y)$ and $(Z,W)$ exchanged. Since $R(X,Y,Z,W)=R(Z,W,X,Y)$, the whole
  computation applies equally with $(X,Y,Z,W)$ replaced by $(Z,W,X,Y)$, giving
  the same left-hand side $(\nabla^*\nabla R)(X,Y,Z,W)$. Averaging the two
  resulting expressions,
  \begin{align*}
    (\nabla^*\nabla R)(X,Y,Z,W)
    &=\tfrac12\big((\nabla_X\nabla^*R)(Y,Z,W)-(\nabla_Y\nabla^*R)(X,Z,W)\big)\\
    &\quad+\tfrac12\big((\nabla_Z\nabla^*R)(W,X,Y)-(\nabla_W\nabla^*R)(Z,X,Y)\big)
      -\tfrac12\operatorname{Ric}(R)(X,Y,Z,W),
  \end{align*}
  which rearranges to the stated identity. Finally, the second Bianchi
  identity shows the right-hand side vanishes identically once $\nabla^*R=0$
  (no additional divergence-of-derivative terms survive, unlike the Hodge
  Laplacian case), so $R$ divergence free implies $\Delta_LR=0$.
\end{proof}

\subsection{Symmetric $(0,2)$-Tensors}

\begin{definition}[the exterior derivative $d^\nabla h$ of a symmetric tensor]
  \label{def:pet-ch9-codazzi-exterior-derivative}
  \lean{PetersenLib.symmetricTensorExteriorDerivative}
  For a symmetric $(0,2)$-tensor $h$ with associated $(1,1)$-tensor $H$
  ($h(z,w)=g(H(z),w)$), recalling $(d^\nabla H)(X,Y)=(\nabla_XH)(Y)-(\nabla_YH)(X)$
  for $(1,1)$-tensors, define
  \[
    d^\nabla h(X,Y,Z)=(\nabla_Xh)(Y,Z)-(\nabla_Yh)(X,Z).
  \]
  \source{petersen:page-0363}
\end{definition}

\begin{remark}[Codazzi identity for the Ricci tensor]
  \label{rem:pet-ch9-ricci-codazzi-identity}
  \lean{PetersenLib.ricciTensor_codazziIdentity}
  \uses{def:pet-ch9-codazzi-exterior-derivative}
  For the second fundamental form $\mathrm{II}$ of an immersed hypersurface
  $M^n\to\mathbb{R}^{n+1}$, the Codazzi--Mainardi equations read $d^\nabla\mathrm{II}=0$.
  For the Ricci tensor, Petersen's exercise 3.4.8 gives
  $(d^\nabla\operatorname{Ric})(X,Y,Z)=(\nabla^*R)(Z,X,Y)$; an analogous identity relates
  the Schouten and Weyl tensors (Petersen, exercise 3.4.26).
  \source{petersen:page-0363}
\end{remark}

\begin{theorem}[Theorem 9.4.3]
  \label{thm:pet-ch9-symmetric-tensor-lichnerowicz-identity}
  \lean{PetersenLib.symmetricTensor_lichnerowiczIdentity}
  \uses{def:pet-ch9-codazzi-exterior-derivative,def:pet-ch9-weitzenbock-curvature-operator}
  Any symmetric $(0,2)$-tensor $h$ on a Riemannian manifold satisfies
  \[
    (\nabla_X\nabla^*h)(X)+(\nabla^*d^\nabla h)(X,X)
    =(\nabla^*\nabla h)(X,X)+\tfrac12(\operatorname{Ric}(h))(X,X).
  \]
  \source{petersen:page-0364}
\end{theorem}
\begin{proof}
  \uses{thm:pet-ch9-symmetric-tensor-lichnerowicz-identity}
  With respect to an orthonormal frame $E_i$ that is normal (parallel) at the
  point of evaluation, the left-hand side terms are
  \[
    (\nabla_X\nabla^*h)(X)=-(\nabla^2_{E_i,E_i}h)(E_i,X)
  \]
  and, expanding $d^\nabla h$ and differentiating,
  \[
    (\nabla^*d^\nabla h)(X,X)=-(\nabla_{E_i}d^\nabla h)(E_i,X,X)
    =-(\nabla^2_{E_i,E_i}h)(X,X)+(\nabla^2_{E_i,X}h)(E_i,X).
  \]
  Adding these,
  \begin{align*}
    (\nabla_X\nabla^*h)(X)+(\nabla^*d^\nabla h)(X,X)
    &=(\nabla^*\nabla h)(X,X)+(\nabla^2_{E_i,X}h)(E_i,X)-(\nabla^2_{E_i,E_i}h)(E_i,X)\\
    &=(\nabla^*\nabla h)(X,X)+\big(R(E_i,X)h\big)(E_i,X),
  \end{align*}
  using the curvature commutation $\nabla^2_{E_i,X}-\nabla^2_{X,E_i}=R(E_i,X)$
  together with $(\nabla^2_{X,E_i}h)(E_i,X)-(\nabla^2_{E_i,E_i}h)(E_i,X)$ folding
  into $(\nabla^*\nabla h)(X,X)$ after resumming over $i$. Since $h$ is
  symmetric, $(R(E_i,X)h)(E_i,X)=\tfrac12(\operatorname{Ric}(h))(X,X)$ (the two terms of
  the Weitzenböck sum coincide by symmetry of $h$ in its two slots), which
  finishes the proof.
\end{proof}

\begin{definition}[Codazzi and harmonic symmetric tensors]
  \label{def:pet-ch9-codazzi-harmonic-tensor}
  \lean{PetersenLib.IsCodazziTensor,PetersenLib.IsHarmonicSymmetricTensor}
  \uses{def:pet-ch9-codazzi-exterior-derivative}
  A symmetric $(0,2)$-tensor $h$ is a \emph{Codazzi tensor} if $d^\nabla h=0$,
  and \emph{harmonic} if in addition it is divergence free ($\nabla^*h=0$).
  \source{petersen:page-0364}
\end{definition}

\begin{proposition}[Proposition 9.4.4]
  \label{prop:pet-ch9-harmonic-iff-codazzi-constant-trace}
  \lean{PetersenLib.harmonic_iff_codazzi_constantTrace}
  \uses{def:pet-ch9-codazzi-harmonic-tensor}
  A symmetric $(0,2)$-tensor is harmonic if and only if it is a Codazzi tensor
  of constant trace.
  \source{petersen:page-0364}
\end{proposition}
\begin{proof}
  \uses{prop:pet-ch9-harmonic-iff-codazzi-constant-trace}
  For any symmetric $(0,2)$-tensor $h$, using a normal frame $E_i$,
  \begin{align*}
    (\nabla^*h)(X)&=-(\nabla_{E_i}h)(E_i,X)=-(\nabla_{E_i}h)(X,E_i)\\
    &=-(\nabla_Xh)(E_i,E_i)+(d^\nabla h)(X,E_i,E_i)\\
    &=-D_X(\operatorname{tr}h)+(d^\nabla h)(X,E_i,E_i).
  \end{align*}
  Hence for a Codazzi tensor ($d^\nabla h=0$), $\nabla^*h=0$ if and only if
  $D_X(\operatorname{tr}h)=0$ for all $X$, i.e.\ $\operatorname{tr}h$ is constant.
\end{proof}

Constant mean curvature hypersurfaces therefore have harmonic second
fundamental form, a fact exploited by both Lichnerowicz and Simons. For the
Ricci tensor, being Codazzi already forces the curvature tensor to be
harmonic, since $\operatorname{Ric}$ is Codazzi exactly when $\nabla^*R=0$
(\cref{rem:pet-ch9-ricci-codazzi-identity}).

\begin{corollary}[Corollary 9.4.5]
  \label{cor:pet-ch9-ricci-harmonic-iff-curvature-harmonic}
  \lean{PetersenLib.ricciTensor_harmonic_iff_curvatureTensor_harmonic}
  \uses{prop:pet-ch9-harmonic-iff-codazzi-constant-trace,rem:pet-ch9-ricci-codazzi-identity,thm:pet-ch9-curvature-tensor-lichnerowicz}
  The Ricci tensor is harmonic if and only if the curvature tensor is
  harmonic (i.e.\ $\nabla^*R=0$).
  \source{petersen:page-0365}
\end{corollary}
\begin{proof}
  \uses{cor:pet-ch9-ricci-harmonic-iff-curvature-harmonic}
  By \cref{rem:pet-ch9-ricci-codazzi-identity}, $\operatorname{Ric}$ is a Codazzi tensor
  exactly when $\nabla^*R=0$. The contracted second Bianchi identity
  (Petersen, proposition 3.1.5) together with the computation in the proof of
  \cref{prop:pet-ch9-harmonic-iff-codazzi-constant-trace} gives
  \[
    2D_X\operatorname{scal}=-(\nabla^*\operatorname{Ric})(X)=D_X(\operatorname{tr}\operatorname{Ric})=D_X(\operatorname{scal}),
  \]
  so $D_X\operatorname{scal}=0$: the scalar curvature is constant, hence
  $\operatorname{Ric}$ is divergence free whenever it is Codazzi. Thus $\operatorname{Ric}$
  is harmonic (Codazzi with constant trace, by
  \cref{prop:pet-ch9-harmonic-iff-codazzi-constant-trace}) if and only if it is
  Codazzi, i.e.\ if and only if $\nabla^*R=0$.
\end{proof}

\subsection{Topological and Geometric Consequences}

\begin{theorem}[Theorem 9.4.6 (D.\ Meyer, 1971, D.\ Meyer--Gallot, 1975, and Gallot, 1981)]
  \label{thm:pet-ch9-nonneg-curvature-operator-parallel-forms}
  \lean{PetersenLib.nonnegCurvatureOperator_harmonicForms_parallel}
  \uses{thm:pet-ch9-weitzenbock-hodge-laplacian,thm:pet-ch9-lichnerowicz-dimension-estimate,cor:pet-ch9-ric-nonneg-curvature-operator}
  Let $(M,g)$ be a closed Riemannian $n$-manifold. If $\mathfrak{R}\ge0$, then
  every harmonic form is parallel. If $\mathfrak{R}>0$, the only parallel
  $l$-forms occur for $l=0,n$. Finally, if $\mathfrak{R}\ge k$ and
  $\operatorname{diam}(M)\le D$,
  \[
    b_l(M)\le\binom{n}{l}\exp\big(D\cdot C(n,kD^2)\sqrt{-kC}\big).
  \]
  \source{petersen:page-0365}
\end{theorem}
\begin{proof}
  \uses{thm:pet-ch9-nonneg-curvature-operator-parallel-forms}
  The first statement follows immediately from
  \cref{thm:pet-ch9-weitzenbock-hodge-laplacian,cor:pet-ch9-ric-nonneg-curvature-operator}:
  when $\mathfrak{R}\ge0$, $g(\operatorname{Ric}(\omega),\omega)\ge0$ for every harmonic
  form $\omega$, so $g(\nabla^*\nabla\omega,\omega)=-g(\operatorname{Ric}(\omega),\omega)\le0$
  and \cref{prop:pet-ch9-vanishing-connection-laplacian} shows $\omega$ is
  parallel (a harmonic form on a closed manifold has $|\omega|^2$ attaining a
  maximum by compactness).

  For the second statement, when $\mathfrak{R}>0$, the identity
  \[
    0=g(\operatorname{Ric}(\omega),\omega)=\sum_\sigma\lambda_\sigma|\Xi_\sigma\omega|^2
  \]
  (from the proof of \cref{cor:pet-ch9-ric-nonneg-curvature-operator}, all
  $\lambda_\sigma>0$) forces $\Xi_\sigma\omega=0$ for every $\sigma$, and by
  linearity $L\omega=0$ for every skew-symmetric $L$. If $0<m<n$ and $L$ is
  chosen with $L(e_i)=0$ for $i<m$ and $L(e_m)=e_{m+1}$, then
  \[
    0=(L\omega)(e_1,\dots,e_m)=-\omega(e_1,\dots,e_{m-1},e_{m+1}),
  \]
  and since the orthonormal basis was arbitrary this forces $\omega=0$ unless
  $l\in\{0,n\}$.

  The last statement is the case $c=1$ of the general estimate
  \cref{thm:pet-ch9-lichnerowicz-dimension-estimate} applied to the
  Lichnerowicz Laplacian $\Delta_L=\square$ on $l$-forms, whose kernel is
  $\mathcal{H}^l(M)$ of dimension $b_l(M)$, with fibre dimension
  $m=\binom{n}{l}$.
\end{proof}

The classification of manifolds with nonnegative curvature operator (Petersen,
theorem 10.3.7) shows there are manifolds with positive or nonnegative
sectional curvature admitting no metric with $\mathfrak{R}\ge0$: $S^2\times S^2$
and $\mathbb{CP}^2$ have $\operatorname{Ric}>0$ but do not admit $\mathfrak{R}\ge0$ (H.\ Hopf's
question for $S^2\times S^2$), and there is a simply connected manifold with
$\operatorname{Ric}>0$ but no metric with $\sec\ge0$ (Chapter 12 of the source).

\begin{example}[Example 9.4.7]
  \label{ex:pet-ch9-cp2-connect-sum-nonneg-curvature}
  \lean{PetersenLib.cp2ConnectSumCp2_nonnegSectionalCurvature}
  \uses{thm:pet-ch9-nonneg-curvature-operator-parallel-forms}
  Let $S^1$ act on $S^3\subset\mathbb{C}^2$ by the Hopf action and on $S^2$ by
  rotation through total angle $2\pi k$; the quotient of $S^2\times S^3$ by
  the diagonal action is $S^2\times S^2$ if $k$ is even and
  $\mathbb{CP}^2\#\mathbb{CP}^2$ if $k$ is odd, and by O'Neill's formula it carries
  nonnegative sectional curvature in either case. By the classification of
  manifolds with $\mathfrak{R}\ge0$ (Petersen, theorem 10.3.7), the only simply
  connected manifolds with $\mathfrak{R}\ge0$ are (topologically) $S^2\times S^2$,
  $S^4$, or $\mathbb{CP}^2$, so this $\mathfrak{R}\ge0$-classification is strictly
  finer than the classification of $\sec\ge0$ manifolds; these examples were
  originally due to Cheeger, via a different construction.
  \source{petersen:page-0366}
\end{example}

The Bochner technique for the Dirac operator's square,
$\nabla^*\nabla+\tfrac14\operatorname{scal}$ (Lichnerowicz, also I.\ Singer), shows the
$\hat A$-genus vanishes on spin manifolds with $\operatorname{scal}>0$; via
generalizations, Gromov--Lawson showed every metric with $\operatorname{scal}\ge0$ on a
torus is flat, generalizing \cref{cor:pet-ch9-betti1-bound-flat-torus}. There
are still no known simply connected examples with $\operatorname{scal}>0$ but no metric
with $\operatorname{Ric}>0$; however for closed $(M,g)$ and small $\varepsilon$,
$(M\times S^2,g+\varepsilon g_{S^2})$ has $\operatorname{scal}>0$, so taking $M$ a surface
with $b_1>4$ gives a non-simply-connected example with $\operatorname{scal}>0$ but
$b_1(M\times S^2)>4$, which by \cref{cor:pet-ch9-betti1-bound-flat-torus}
cannot support $\operatorname{Ric}\ge0$.

\begin{theorem}[Theorem 9.4.8 (Tachibana, 1974)]
  \label{thm:pet-ch9-tachibana-parallel-curvature}
  \lean{PetersenLib.tachibana_parallelCurvature}
  \uses{thm:pet-ch9-curvature-tensor-lichnerowicz,prop:pet-ch9-vanishing-connection-laplacian,conv:pet-ch9-so-metric-convention}
  Let $(M,g)$ be closed. If $\mathfrak{R}\ge0$ and $\nabla^*R=0$, then $\nabla R=0$.
  If moreover $\mathfrak{R}>0$, then $(M,g)$ has constant curvature.
  \source{petersen:page-0367}
\end{theorem}
\begin{proof}
  \uses{thm:pet-ch9-tachibana-parallel-curvature}
  By \cref{thm:pet-ch9-curvature-tensor-lichnerowicz}, $\nabla^*R=0$ gives
  $\nabla^*\nabla R+\tfrac12\operatorname{Ric}(R)=0$. If $\mathfrak{R}\ge0$, then as in
  \cref{cor:pet-ch9-ric-nonneg-curvature-operator}, $g(\operatorname{Ric}(R),R)\ge0$, so
  $g(\nabla^*\nabla R,R)\le0$; since $M$ is closed, $|R|^2$ has a maximum, and
  \cref{prop:pet-ch9-vanishing-connection-laplacian} gives $\nabla R=0$.

  If moreover $\mathfrak{R}>0$, as in the proof of
  \cref{thm:pet-ch9-nonneg-curvature-operator-parallel-forms} one gets $LR=0$
  for every $L\in\mathfrak{so}(T_pM)$. To show the mixed curvatures
  $R(x,v,w,z)$ vanish for $z$ perpendicular to $x,v,w$: choose $L$ with
  $L(y)=0$ and $L(x)=z$; then
  \[
    0=LR(x,y,y,x)=-R(L(x),y,y,x)-R(x,y,y,L(x))=-2R(x,y,y,z).
  \]
  Polarizing in $y=v+w$ gives $R(x,v,w,z)=-R(x,w,v,z)$. The (first) Bianchi
  identity then gives
  \[
    R(x,v,w,z)=R(w,x,z,v)-R(w,x,v,z)=-2R(w,x,v,z)=2R(x,w,v,z)=-2R(x,v,w,z),
  \]
  forcing $R(x,v,w,z)=0$. Consequently every bivector $x\wedge y$ is an
  eigenvector of $\mathfrak{R}$ (all mixed curvatures with a fourth
  perpendicular direction vanish, and $R(x,y,y,x)$ alone survives), which
  forces $\mathfrak{R}=kI$ for a constant $k$, i.e.\ $(M,g)$ has constant
  curvature.
\end{proof}

\subsection{Simplification of $g(\operatorname{Ric}(T),T)$}

An alternate method recovers the $1$-form formula and extends to general
$(0,2)$-tensors, replacing the orthonormal basis $\Xi_\sigma$ by a single
$\Lambda^2TM$-valued tensor built from $T$.

\begin{definition}[$\Lambda^2TM$-valued tensor $\hat T$]
  \label{def:pet-ch9-tensor-lambda2-valued}
  \lean{PetersenLib.lambda2ValuedTensor}
  \uses{conv:pet-ch9-so-metric-convention}
  A $(0,k)$-tensor $T$ determines $\hat T$, valued in $\Lambda^2TM\cong
  \mathfrak{so}(TM)$, implicitly by
  \[
    g\big(L,\hat T(X_1,\dots,X_k)\big)=(LT)(X_1,\dots,X_k)\quad\text{for all }
    L\in\mathfrak{so}(TM).
  \]
  \source{petersen:page-0368}
\end{definition}

\begin{lemma}[Lemma 9.4.9]
  \label{lem:pet-ch9-ric-curvature-operator-hat-formula}
  \lean{PetersenLib.weitzenbock_eq_curvatureOperator_hat}
  \uses{def:pet-ch9-tensor-lambda2-valued,def:pet-ch9-weitzenbock-curvature-operator,lem:pet-ch9-ric-curvature-operator-formula}
  For all $(0,k)$-tensors $T,S$, $g(\operatorname{Ric}(T),S)=g(\mathfrak{R}(\hat T),\hat S)$.
  \source{petersen:page-0368}
\end{lemma}
\begin{proof}
  \uses{lem:pet-ch9-ric-curvature-operator-hat-formula}
  Using \cref{lem:pet-ch9-ric-curvature-operator-formula} and summing over a
  multi-index $I=(i_1,\dots,i_k)$ and an orthonormal basis $\Xi_\sigma$ of
  $\mathfrak{so}(T_pM)$,
  \begin{align*}
    g(\operatorname{Ric}(T),S)&=\sum_\sigma g(\Xi_\sigma T,R(\Xi_\sigma)S)\\
    &=\sum_{\sigma,I}(\Xi_\sigma T)(e_I)\,(R(\Xi_\sigma)S)(e_I)\\
    &=\sum_{\sigma,I}g\big(\Xi_\sigma,\hat T(e_I)\big)\,g\big(R(\Xi_\sigma),\hat S(e_I)\big)\\
    &=\sum_I g\Big(\mathfrak{R}\big(\hat T(e_I)\big),\hat S(e_I)\Big)\\
    &=g(\mathfrak{R}(\hat T),\hat S),
  \end{align*}
  where the third equality uses the defining identity of
  $\hat T,\hat S$ from \cref{def:pet-ch9-tensor-lambda2-valued}, and the fourth
  expands $R(\Xi_\sigma)=\sum_\tau g(\mathfrak{R}(\Xi_\sigma),\Xi_\tau)\Xi_\tau$ and
  resums over $\sigma$ using that $\Xi_\sigma$ is an orthonormal basis. Taking
  $S=T$ recovers self-adjointness of $\operatorname{Ric}$, now manifest since
  $\mathfrak{R}$ is self-adjoint on $\Lambda^2TM$.
\end{proof}

This reformulation of $g(\operatorname{Ric}(T),T)$ is again manifestly nonnegative when
$\mathfrak{R}\ge0$, and it sometimes shows nonnegativity under weaker
hypotheses, as worked out below for $1$-forms and symmetric $(0,2)$-tensors.

\begin{proposition}[Proposition 9.4.10]
  \label{prop:pet-ch9-1form-hat-ricci}
  \lean{PetersenLib.oneForm_hat_ricci_identity}
  \uses{lem:pet-ch9-ric-curvature-operator-hat-formula,def:pet-ch9-tensor-lambda2-valued}
  If $\omega$ is a $1$-form dual to the vector field $X$, then
  $\hat\omega(Z)=X\wedge Z$ and $g(\mathfrak{R}(\hat\omega),\hat\omega)=\operatorname{Ric}(X,X)$.
  \source{petersen:page-0368,petersen:page-0369}
\end{proposition}
\begin{proof}
  \uses{prop:pet-ch9-1form-hat-ricci}
  For $L\in\mathfrak{so}(T_pM)$, $(L\omega)(Z)=-\omega(L(Z))=-g(X,L(Z))
  =-g(L,Z\wedge X)=g(L,X\wedge Z)$, so by the defining property of
  $\hat\omega$, $\hat\omega(Z)=X\wedge Z$. Then, using
  \cref{lem:pet-ch9-ric-curvature-operator-hat-formula} with $S=T=\omega$ and
  an orthonormal frame $E_i$,
  \begin{align*}
    g(\operatorname{Ric}(\omega),\omega)
    &=g(\mathfrak{R}(\hat\omega),\hat\omega)
    =\sum_ig\big(\hat\omega(E_i),\mathfrak{R}(\hat\omega(E_i))\big)\\
    &=\sum_ig\big(\mathfrak{R}(E_i\wedge X),E_i\wedge X\big)
    =\sum_iR(X,E_i,E_i,X)=\operatorname{Ric}(X,X).\qedhere
  \end{align*}
\end{proof}

\begin{remark}[$p$-form generalization]
  \label{rem:pet-ch9-pform-hat-formula}
  \lean{PetersenLib.pForm_hat_formula}
  \uses{def:pet-ch9-tensor-lambda2-valued,prop:pet-ch9-1form-hat-ricci}
  More generally, if $\omega$ is a $p$-form and $\Omega$ is the
  $(p-1)$-form-valued map with $g(\Omega(X_1,\dots,X_{p-1}),X_p)=
  \omega(X_1,\dots,X_p)$, then
  \[
    \hat\omega(X_1,\dots,X_p)=\sum_{i=1}^p(-1)^{p-i}X_i\wedge
    \Omega(X_1,\dots,\hat X_i,\dots,X_p),
  \]
  and $\hat\omega$ vanishes only where $\omega$ does.
  \source{petersen:page-0369}
\end{remark}

Turning to $(0,2)$-tensors, let $h(z,w)=g(H(z),w)$ with adjoint $H^*$.

\begin{proposition}[Proposition 9.4.11]
  \label{prop:pet-ch9-02tensor-hat-formula}
  \lean{PetersenLib.symmetricTensor_hat_formula}
  \uses{def:pet-ch9-tensor-lambda2-valued}
  With $h,H,H^*$ as above, $\hat h(z,w)=H(z)\wedge w-z\wedge H^*(w)$, and
  $\hat h=0$ if and only if $h=\lambda g$ for a scalar $\lambda$.
  \source{petersen:page-0369,petersen:page-0370}
\end{proposition}
\begin{proof}
  \uses{prop:pet-ch9-02tensor-hat-formula}
  For $L\in\mathfrak{so}(T_pM)$,
  \begin{align*}
    (Lh)(z,w)&=-h(L(z),w)-h(z,L(w))\\
    &=-g(H(L(z)),w)-g(H(z),L(w))\\
    &=-g(L(z),H^*(w))-g(L(w),H(z))\\
    &=-g(L,z\wedge H^*(w))-g(L,w\wedge H(z))\\
    &=g\big(L,H(z)\wedge w-z\wedge H^*(w)\big),
  \end{align*}
  which gives the stated formula for $\hat h$. If $h=\lambda g$ then
  $H=H^*=\lambda I$, so $\hat h=0$. Conversely suppose $\hat h=0$; then for all
  $z,w$,
  \[
    z\wedge H^*(H(w))=H(z)\wedge H(w)=-H(w)\wedge H(z)=-w\wedge H^*(H(z))
    =H^*(H(z))\wedge w.
  \]
  This forces $H^*H=\lambda^2I$ and $H=\lambda I$ for a scalar $\lambda$: if
  $H$ were not a multiple of the identity, choosing $z,w$ as distinct
  eigenvectors (or a generic pair) of the self-adjoint $H^*H$ makes the two
  bivectors $z\wedge H^*(H(w))$ and $H^*(H(z))\wedge w$ fail to coincide unless
  $H^*H$ acts as a scalar and $H(z)$ is proportional to $z$ for every $z$,
  i.e.\ $H=\lambda I$; then $H^*=\lambda I=H$, so $h=\lambda g$.
\end{proof}

If $H$ is normal it diagonalizes over $T_pM\otimes\mathbb{C}$: say $H(z)=\lambda z$,
$H(w)=\mu w$ for orthonormal $z,w\in T_pM\otimes\mathbb{C}$. Then
\[
  g\big(\mathfrak{R}(H(z)\wedge w-z\wedge H^*(w)),H(z)\wedge w-z\wedge H^*(w)\big)
  =|\lambda-\bar\mu|^2\,g\big(\mathfrak{R}(\bar z\wedge w),\bar z\wedge w\big).
\]

\begin{definition}[complex sectional curvature]
  \label{def:pet-ch9-complex-sectional-curvature}
  \lean{PetersenLib.complexSectionalCurvature}
  \uses{prop:pet-ch9-02tensor-hat-formula}
  For $z,w\in T_pM\otimes\mathbb{C}$, the \emph{complex sectional curvature} is
  $g(\mathfrak{R}(z\wedge w),\bar z\wedge\bar w)$.
  \source{petersen:page-0370}
\end{definition}

\begin{lemma}[real expansion of the complex sectional curvature]
  \label{lem:pet-ch9-complex-sectional-curvature-real-formula}
  \lean{PetersenLib.complexSectionalCurvature_realFormula}
  \uses{def:pet-ch9-complex-sectional-curvature}
  Writing $z=x+iy$, $w=u+iv$ with $x,y,u,v\in T_pM$,
  \begin{align*}
    g(\mathfrak{R}(z\wedge w),\bar z\wedge\bar w)
    &=R(x,u,u,x)+R(y,v,v,y)+R(x,v,v,x)+R(y,u,u,y)+2R(y,x,v,u).
  \end{align*}
  In particular this quantity is nonnegative whenever $\mathfrak{R}\ge0$, so
  nonnegative complex sectional curvature is a weaker hypothesis than
  $\mathfrak{R}\ge0$ but a stronger one than nonnegative sectional curvature.
  \source{petersen:page-0370,petersen:page-0371}
\end{lemma}
\begin{proof}
  \uses{lem:pet-ch9-complex-sectional-curvature-real-formula}
  Expanding $z\wedge w=(x+iy)\wedge(u+iv)=(x\wedge u-y\wedge v)+i(x\wedge v+y\wedge u)$
  and $\bar z\wedge\bar w=(x\wedge u-y\wedge v)-i(x\wedge v+y\wedge u)$, bilinearity
  and $\mathbb{C}$-linearity of $\mathfrak{R}$ (extended complex-bilinearly) give
  \begin{align*}
    g(\mathfrak{R}(z\wedge w),\bar z\wedge\bar w)
    &=g\big(\mathfrak{R}(x\wedge u-y\wedge v),x\wedge u-y\wedge v\big)
      +g\big(\mathfrak{R}(x\wedge v+y\wedge u),x\wedge v+y\wedge u\big)\\
    &=g(\mathfrak{R}(x\wedge u),x\wedge u)+g(\mathfrak{R}(y\wedge v),y\wedge v)
      +g(\mathfrak{R}(x\wedge v),x\wedge v)+g(\mathfrak{R}(y\wedge u),y\wedge u)\\
    &\quad-2g(\mathfrak{R}(x\wedge u),y\wedge v)+2g(\mathfrak{R}(x\wedge v),y\wedge u)\\
    &=R(x,u,u,x)+R(y,v,v,y)+R(x,v,v,x)+R(y,u,u,y)\\
    &\quad+2R(x,u,y,v)-2R(x,v,y,u)\\
    &=R(x,u,u,x)+R(y,v,v,y)+R(x,v,v,x)+R(y,u,u,y)-2\big(R(v,y,x,u)+R(x,v,y,u)\big)\\
    &=R(x,u,u,x)+R(y,v,v,y)+R(x,v,v,x)+R(y,u,u,y)+2R(y,x,v,u),
  \end{align*}
  where the last two lines use the first Bianchi identity to combine the two
  mixed curvature terms into $2R(y,x,v,u)$. Setting $y=v=0$ this reads
  $g(\mathfrak{R}(z\wedge w),\bar z\wedge\bar w)=R(x,u,u,x)$, the ordinary
  sectional-curvature numerator, showing the first line already exhibits the
  quantity as a nonnegative combination of sectional-curvature terms plus a
  correction whenever $\mathfrak{R}\ge0$.
\end{proof}

\begin{remark}[special cases of the complex sectional curvature]
  \label{rem:pet-ch9-complex-sectional-curvature-special-cases}
  \lean{PetersenLib.complexSectionalCurvature_specialCases}
  \uses{lem:pet-ch9-complex-sectional-curvature-real-formula}
  Depending on $\dim_{\mathbb{R}}\operatorname{span}\{x,y,u,v\}$: if $y=v=0$ the complex
  sectional curvature is the ordinary sectional curvature $R(x,u,u,x)$; if
  $x,y,u,v$ are orthonormal it is called the \emph{isotropic curvature}; and if
  $u=v$ it reduces to $2R(x,u,u,x)+2R(y,u,u,y)$, called the \emph{second Ricci
  curvature} when $x,y,u$ are orthonormal.
  \source{petersen:page-0371}
\end{remark}

\begin{proposition}[Proposition 9.4.12]
  \label{prop:pet-ch9-normal-tensor-nonneg-complex-sectional}
  \lean{PetersenLib.normalTensor_nonneg_of_complexSectionalCurvature_nonneg}
  \uses{prop:pet-ch9-02tensor-hat-formula,def:pet-ch9-complex-sectional-curvature}
  Let $h$ be a $(0,2)$-tensor with $H$ normal. If all complex sectional
  curvatures are nonnegative, then $g(\mathfrak{R}(\hat h),\hat h)\ge0$.
  \source{petersen:page-0371}
\end{proposition}
\begin{proof}
  \uses{prop:pet-ch9-normal-tensor-nonneg-complex-sectional}
  Since $H$ is normal it has a complex orthonormal eigenbasis $e_i$,
  $H(e_i)=\lambda_ie_i$, $H^*(e_i)=\bar\lambda_ie_i$. By
  \cref{prop:pet-ch9-02tensor-hat-formula},
  \begin{align*}
    g(\mathfrak{R}(\hat h),\hat h)
    &=\sum_{i,j}g\big(\mathfrak{R}(\hat h(e_i,e_j)),\hat h(e_i,e_j)\big)\\
    &=\sum_{i,j}g\big(\mathfrak{R}(H(e_i)\wedge e_j-e_i\wedge H^*(e_j)),
       H(e_i)\wedge e_j-e_i\wedge H^*(e_j)\big)\\
    &=\sum_{i,j}|\lambda_i-\bar\lambda_j|^2\,g\big(\mathfrak{R}(e_i\wedge e_j),\bar e_i\wedge\bar e_j\big),
  \end{align*}
  a nonnegative combination of complex sectional curvatures, hence
  nonnegative. When $H$ is self-adjoint the $\lambda_i$ are real and only
  ordinary (real) sectional curvatures appear; when $H$ is skew-adjoint its
  nonzero eigenvalues are purely imaginary, so isotropic curvatures and (in
  odd dimensions, from the zero eigenvalue) second Ricci curvatures appear
  instead, with no real sectional curvature terms.
\end{proof}

\begin{proposition}[Proposition 9.4.13]
  \label{prop:pet-ch9-general-tensor-nonneg-complex-sectional}
  \lean{PetersenLib.tensor_nonneg_of_complexSectionalCurvature_nonneg}
  \uses{prop:pet-ch9-normal-tensor-nonneg-complex-sectional}
  $g(\mathfrak{R}(\hat h),\hat h)\ge0$ for every $(0,2)$-tensor $h$ on $T_pM$, as
  soon as every complex sectional curvature on $T_pM$ is nonnegative.
  \source{petersen:page-0372}
\end{proposition}
\begin{proof}
  \uses{prop:pet-ch9-general-tensor-nonneg-complex-sectional}
  Write $h=h_s+h_a$ for the symmetric and skew-symmetric parts (both normal
  operators). Bilinearity of $g(\mathfrak{R}(-),-)$ gives
  \[
    g(\mathfrak{R}(\hat h),\hat h)=g(\mathfrak{R}(\hat h_s),\hat h_s)
    +g(\mathfrak{R}(\hat h_a),\hat h_a)+2g(\mathfrak{R}(\hat h_s),\hat h_a),
  \]
  using symmetry of $g(\mathfrak{R}(-),-)$ for the cross term. But, using an
  orthonormal frame $e_i$ and antisymmetry of $\hat h_a(e_i,e_j)$ in $i,j$
  together with symmetry of $\hat h_s(e_i,e_j)$,
  \[
    g(\mathfrak{R}(\hat h_s),\hat h_a)=\sum_{i,j}g\big(\mathfrak{R}(\hat h_s(e_i,e_j)),\hat h_a(e_i,e_j)\big)
    =-\sum_{i,j}g\big(\mathfrak{R}(\hat h_s(e_j,e_i)),\hat h_a(e_j,e_i)\big)
    =-g(\mathfrak{R}(\hat h_s),\hat h_a),
  \]
  so this cross term vanishes. Hence
  $g(\mathfrak{R}(\hat h),\hat h)=g(\mathfrak{R}(\hat h_s),\hat h_s)+g(\mathfrak{R}(\hat h_a),\hat h_a)$,
  and each summand is nonnegative by
  \cref{prop:pet-ch9-normal-tensor-nonneg-complex-sectional} applied to $h_s$
  and to $h_a$ (both normal, since symmetric and skew-symmetric operators are
  normal).
\end{proof}

\section{Further Study}

For general and complete accounts of the Bochner technique and spin geometry,
including a complete proof of the Hodge theorem
(\cref{thm:pet-ch9-hodge-theorem}), the source recommends two texts (labelled
[107] and [71] in its bibliography), with further treatments of the Hodge
theorem in [65], [92], and [101]. For generalizations to manifolds with
integral curvature bounds, extending the Betti-number estimates of
\cref{thm:pet-ch9-gromov-gallot-betti1-estimate,thm:pet-ch9-nonneg-curvature-operator-parallel-forms},
the source points to [46].

\section{Exercises}

\begin{remark}[Exercise 9.6.1]
  \label{rem:pet-ch9-exercise-9-6-1}
  \lean{PetersenLib.exercise9_6_1}
  \uses{cor:pet-ch9-betti1-bound-flat-torus,thm:pet-ch9-bochner-vanishing-1forms}
  Suppose $(M^n,g)$ is compact with $b_1=k$. If $\operatorname{Ric}\ge0$, show that the
  universal cover splits, $(\tilde M,g)=(N,h)\times(\mathbb{R}^k,g_{\mathbb{R}^k})$. Give an
  example where $b_1<n$ and $(\tilde M,g)=(\mathbb{R}^n,g_{\mathbb{R}^n})$.
  \source{petersen:page-0373}
\end{remark}

\begin{remark}[Exercise 9.6.2]
  \label{rem:pet-ch9-exercise-9-6-2}
  \lean{PetersenLib.exercise9_6_2}
  \uses{prop:pet-ch9-vector-field-1form-duality}
  Show directly that if $E_i$ is an orthonormal frame and $v\mapsto\nabla_vX$
  is symmetric, then $\sum_ig(\nabla_{E_i}X,\nabla_XE_i)=0$, without assuming
  the $E_i$ are parallel along $X$.
  \source{petersen:page-0373}
\end{remark}

\begin{remark}[Exercise 9.6.3]
  \label{rem:pet-ch9-exercise-9-6-3}
  \lean{PetersenLib.exercise9_6_3}
  \uses{def:pet-ch9-codifferential}
  Show that for an oriented Riemannian manifold,
  $(\operatorname{im}d^{k-1})^\perp=\ker\delta^{k-1}$ and $\operatorname{im}d^{k-1}\subset
  (\ker\delta^{k-1})^\perp$ in $\Omega^k(M)$.
  \source{petersen:page-0373}
\end{remark}

\begin{remark}[Exercise 9.6.4]
  \label{rem:pet-ch9-exercise-9-6-4}
  \lean{PetersenLib.exercise9_6_4}
  \uses{def:pet-ch9-lichnerowicz-laplacian}
  Let $(M^n,g)$ be compact and $E\to M$ a tensor bundle. Let $W^{k,2}(E)$ be the
  Hilbert space completion of $\Gamma(E)$ with square norm
  $\sum_{i=0}^k\|\nabla^iT\|_2^2$, so that $T\in W^{1,2}(E)$ has an $L^2$
  derivative $\nabla T$ satisfying $(\nabla T,\nabla S)=(T,\nabla^*\nabla S)$ for
  all $S\in\Gamma(E)$; $\Gamma(E)\subset W^{k,2}(E)$ is dense, the Sobolev
  inequality gives $\bigcap_kW^{k,2}(E)=\bigcap_kW^{k,p}(E)$ for all $p$, and
  the techniques of section 7.1.5 show a tensor in $W^{1,p}$ is H\"older
  continuous for $p>n$, so $\Gamma(E)=\bigcap_kW^{k,2}(E)$.
  \begin{enumerate}
    \item[(1)] Show that for all $T\in\Gamma(E)$ there is a commutation
      relation $\nabla^*\nabla(\nabla^kT)-\nabla^k(\nabla^*\nabla T)
      =\sum_{i=0}^kC_i^k(\nabla^{k-1}R\otimes\nabla^iT)$ for suitable
      contractions $C_i^k$.
    \item[(2)] If $T\in W^{1,2}(E)$ and $T'\in W^{k,2}(E)$ satisfy
      $(T',S)=(\nabla T,\nabla S)$ for all $S\in\Gamma(E)$ (i.e.\
      $\nabla^*\nabla T=T'$ weakly), show $T\in W^{k+1,2}(E)$, by defining weak
      derivatives $\nabla^{l+1}T$ inductively via
      $(\nabla^{l+1}T,\nabla^{l+1}S)=(\nabla^lT',\nabla^lS)
      +\sum_{i=0}^l(C_i^l(\nabla^{l-i}R\otimes\nabla^iT),\nabla^lS)$.
    \item[(3)] Conclude that if $T'\in\Gamma(E)$, then $T\in\Gamma(E)$ and
      $\nabla^*\nabla T=T'$.
  \end{enumerate}
  \source{petersen:page-0373}
\end{remark}

\begin{remark}[Exercise 9.6.5]
  \label{rem:pet-ch9-exercise-9-6-5}
  \lean{PetersenLib.exercise9_6_5}
  \uses{def:pet-ch9-lichnerowicz-laplacian,rem:pet-ch9-exercise-9-6-4,cor:pet-ch9-ric-nonneg-curvature-operator}
  Let $(M^n,g)$ be compact with $\operatorname{diam}M\le D$, $\operatorname{Ric}\ge k$, and $E\to M$ a
  tensor bundle of fibre dimension $m$ with Lichnerowicz Laplacian $\Delta_L$.
  Establish the spectral theorem for $\Delta_L$ and conclude the orthogonal
  decomposition $\Gamma(E)=\ker\Delta_L\oplus\operatorname{im}\Delta_L$.
  \begin{enumerate}
    \item[(1)] Using the completion $W^{1,2}(E)$ of Exercise 9.6.4, show that
      $(\Delta_LT,S)=(\nabla T,\nabla S)+c(\operatorname{Ric}(T),S)$ is symmetric and
      well-defined for $T,S\in W^{1,2}(E)$.
    \item[(2)] Show $\inf_{T\in W^{1,2},\|T\|_2=1}(\Delta_LT,T)>cCk$ for $k\le0$,
      with $C$ the constant of Corollary 9.3.4.
    \item[(3)] Show a sequence $T_i\subset W^{1,2}(E)$ with $\|T_i\|_2=1$ and
      $(\Delta_LT_i,T_i)$ bounded has an $L^2$-convergent subsequence that is
      also weakly convergent in $W^{1,2}(E)$ (use theorem 7.1.18).
    \item[(4)] For a closed $\Delta_L$-invariant subspace $V\subset W^{1,2}(E)$,
      show the infimum $\lambda=\inf_{T\in V,\|T\|_2=1}(\Delta_LT,T)$ is
      achieved by some $T\in V$, and using (1)–(3) of Exercise 9.6.4 that
      $T\in\Gamma(E)$ and $\Delta_LT=\lambda T$.
    \item[(5)] For a finite-dimensional $V\subset\Gamma(E)$ spanned by
      eigentensors $\Delta_LT_i=\lambda_iT_i$, show
      $\dim V\le mC(n,\max_i\lambda_i,c,D^2k)$.
    \item[(6)] Show all eigenspaces of $\Delta_L$ are finite dimensional and
      the eigenvalues form a discrete set $\lambda_1<\lambda_2<\cdots\to\infty$.
    \item[(7)] Show the eigenspaces are orthogonal with dense direct sum in
      $\Gamma(E)$.
    \item[(8)] Show $\Gamma(E)=\ker\Delta_L\oplus\operatorname{im}\Delta_L$.
  \end{enumerate}
  \source{petersen:page-0373,petersen:page-0374}
\end{remark}

\begin{remark}[Exercise 9.6.6]
  \label{rem:pet-ch9-exercise-9-6-6}
  \lean{PetersenLib.exercise9_6_6}
  \uses{thm:pet-ch9-bochner-vanishing-1forms}
  Let $(M,g)$ be $n$-dimensional and isometric to Euclidean space outside a
  compact $K\subset M$, i.e.\ $M-K$ is isometric to $\mathbb{R}^n-C$ for compact
  $C\subset\mathbb{R}^n$. If $\operatorname{Ric}_g\ge0$, show $M=\mathbb{R}^n$. (Hint: find a metric on
  the $n$-torus isometric to a neighborhood of $K$ and otherwise flat, or show
  a parallel $1$-form on $\mathbb{R}^n-C$ extends to a harmonic $1$-form on $M$ that
  Bochner's formula forces to be parallel.)
  \source{petersen:page-0374}
\end{remark}

\begin{remark}[Exercise 9.6.7]
  \label{rem:pet-ch9-exercise-9-6-7}
  \lean{PetersenLib.exercise9_6_7}
  \uses{def:pet-ch9-lichnerowicz-laplacian,thm:pet-ch9-weitzenbock-hodge-laplacian}
  Let $(M,g)$ be Einstein. Show that all harmonic $1$-forms are eigenforms for
  the connection Laplacian $\nabla^*\nabla$.
  \source{petersen:page-0374}
\end{remark}

\begin{remark}[Exercise 9.6.8]
  \label{rem:pet-ch9-exercise-9-6-8}
  \lean{PetersenLib.exercise9_6_8}
  \uses{prop:pet-ch9-bochner-formula-1forms}
  Given vector fields $X,Y$ on $(M,g)$ with $\nabla X,\nabla Y$ symmetric,
  develop Bochner formulas for $\operatorname{Hess}\tfrac12g(X,Y)$ and
  $\Delta\tfrac12g(X,Y)$.
  \source{petersen:page-0374}
\end{remark}

\begin{remark}[Exercise 9.6.9]
  \label{rem:pet-ch9-exercise-9-6-9}
  \lean{PetersenLib.exercise9_6_9}
  \uses{prop:pet-ch9-vanishing-connection-laplacian}
  For tensors $s_1,s_2$ of the same type, show, in analogy with
  $\Delta\tfrac12|s|^2=|\nabla s|^2-g(\nabla^*\nabla s,s)$, that
  \[
    \Delta g(s_1,s_2)=2g(\nabla s_1,\nabla s_2)+g(\nabla^*\nabla s_1,s_2)
    +g(s_1,\nabla^*\nabla s_2).
  \]
  Use this on forms to develop Bochner formulas for inner products of such
  sections. More generally, for the $1$-form $\omega(v)=g(\nabla_vs_1,s_2)$
  (representing half the differential of $g(s_1,s_2)$), show
  \[
    -\delta\omega=g(\nabla s_1,\nabla s_2)-g(\nabla^*\nabla s_1,s_2),
  \]
  \[
    d\omega(X,Y)=g(R(X,Y)s_1,s_2)-g(\nabla_Xs_1,\nabla_Ys_2)+g(\nabla_Ys_1,\nabla_Xs_2).
  \]
  \source{petersen:page-0374,petersen:page-0375}
\end{remark}

\begin{remark}[Exercise 9.6.10]
  \label{rem:pet-ch9-exercise-9-6-10}
  \lean{PetersenLib.exercise9_6_10}
  \uses{conv:pet-ch9-so-metric-convention}
  Let $(M,g)$ be $n$-dimensional.
  \begin{enumerate}
    \item[(1)] Show $L\omega=0$ if $L$ is skew-symmetric and $\omega$ is an
      $n$-form.
    \item[(2)] When $n=2$, show $LR=0$ for all skew-symmetric $L$.
    \item[(3)] For general $L$, show $L\,\mathrm{vol}=\operatorname{tr}(L)\,\mathrm{vol}$.
  \end{enumerate}
  \source{petersen:page-0375}
\end{remark}

\begin{remark}[Exercise 9.6.11 (Simons)]
  \label{rem:pet-ch9-exercise-9-6-11}
  \lean{PetersenLib.exercise9_6_11}
  \uses{def:pet-ch9-codazzi-harmonic-tensor,thm:pet-ch9-symmetric-tensor-lichnerowicz-identity,prop:pet-ch9-vanishing-connection-laplacian}
  Let $(M,g)$ be compact with a symmetric Codazzi tensor $h$ of constant
  trace.
  \begin{enumerate}
    \item[(1)] If $\sec\ge0$, show $\nabla h=0$.
    \item[(2)] If $\sec>0$, show $h=c\cdot g$ for a constant $c$.
    \item[(3)] If the Gauss equations $R(X,Y,Z,W)=h(X,W)h(Y,Z)-h(X,Z)h(Y,W)$
      hold and $\operatorname{tr}h=0$, show
      $\Delta\tfrac12|h|^2\ge|\nabla h|^2-|h|^4$.
  \end{enumerate}
  \source{petersen:page-0375}
\end{remark}

\begin{remark}[Exercise 9.6.12]
  \label{rem:pet-ch9-exercise-9-6-12}
  \lean{PetersenLib.exercise9_6_12}
  \uses{def:pet-ch9-codazzi-exterior-derivative,rem:pet-ch9-exercise-9-6-11}
  Let $(M^n,g)\hookrightarrow\mathbb{R}^{n+1}$ be an isometric immersion.
  \begin{enumerate}
    \item[(1)] Show the second fundamental form $\mathrm{II}$ is a Codazzi
      tensor.
    \item[(2)] Show Liebmann's theorem: if $(M,g)$ has constant mean curvature
      and nonnegative second fundamental form, then $(M,g)$ is a constant
      curvature sphere. (Wente exhibits immersed tori of constant mean
      curvature, so some positivity hypothesis is necessary.)
  \end{enumerate}
  \source{petersen:page-0375,petersen:page-0376}
\end{remark}

\begin{remark}[Exercise 9.6.13 (Berger)]
  \label{rem:pet-ch9-exercise-9-6-13}
  \lean{PetersenLib.exercise9_6_13}
  \uses{cor:pet-ch9-ricci-harmonic-iff-curvature-harmonic,thm:pet-ch9-symmetric-tensor-lichnerowicz-identity}
  Show that a compact manifold with harmonic curvature and nonnegative
  sectional curvature has parallel Ricci curvature.
  \source{petersen:page-0376}
\end{remark}

\begin{remark}[Exercise 9.6.14]
  \label{rem:pet-ch9-exercise-9-6-14}
  \lean{PetersenLib.exercise9_6_14}
  \uses{thm:pet-ch9-hodge-theorem,def:pet-ch9-hodge-cohomology}
  Let $K$ be a Killing field on a closed $(M,g)$ and $\omega$ a harmonic form.
  \begin{enumerate}
    \item[(1)] Show $\mathcal{L}_K\omega=0$ (hint: show $\mathcal{L}_K\omega$
      is also harmonic).
    \item[(2)] Show $i_K\omega$ is closed but not necessarily harmonic.
    \item[(3)] Show all harmonic forms are invariant under $\operatorname{Iso}_0(M)$.
    \item[(4)] Give an example where a harmonic form is not invariant under
      all of $\operatorname{Iso}(M)$.
  \end{enumerate}
  \source{petersen:page-0376}
\end{remark}

\begin{remark}[Exercise 9.6.15]
  \label{rem:pet-ch9-exercise-9-6-15}
  \lean{PetersenLib.exercise9_6_15}
  \uses{rem:pet-ch9-poincare-duality}
  Let $(M,g)$ be a closed K\"ahler manifold with K\"ahler form $\omega$ (a
  parallel nondegenerate $2$-form). Show $\omega^k=\omega\wedge\cdots\wedge
  \omega$ ($k$ times) is closed but not exact, by showing
  $\omega^{\dim M/2}$ is proportional to the volume form; conclude that no
  even homology group vanishes.
  \source{petersen:page-0376}
\end{remark}

\begin{remark}[Exercise 9.6.16]
  \label{rem:pet-ch9-exercise-9-6-16}
  \lean{PetersenLib.exercise9_6_16}
  \uses{def:pet-ch9-codazzi-exterior-derivative,def:pet-ch9-weitzenbock-curvature-operator}
  Let $E\to M$ be a tensor bundle.
  \begin{enumerate}
    \item[(1)] Let $\Omega^p(M,E)$ be the alternating $p$-linear maps
      $TM\to E$ (so $\Omega^0(M,E)=\Gamma(E)$). Show $\Omega^*(M)$ acts
      naturally on both sides of $\Omega^*(M,E)$ by wedge product.
    \item[(2)] Show there is a natural wedge product
      $\Omega^p(M,\operatorname{Hom}(E,E))\times\Omega^q(M,E)\to\Omega^{p+q}(M,E)$.
    \item[(3)] Show there is a connection-dependent exterior derivative
      $d^\nabla:\Omega^p(M,E)\to\Omega^{p+1}(M,E)$ satisfying the Leibniz rule
      for the above wedge products, with $d^\nabla s=\nabla s$ for
      $s\in\Gamma(E)$.
    \item[(4)] Viewing $R(X,Y)s\in\Omega^2(M,\operatorname{Hom}(E,E))$, show
      $(d^\nabla\circ d^\nabla)(s)=R\wedge s$ for $s\in\Omega^p(M,E)$, and that
      the second Bianchi identity reads $d^\nabla R=0$.
  \end{enumerate}
  \source{petersen:page-0376}
\end{remark}

\begin{remark}[Exercise 9.6.17]
  \label{rem:pet-ch9-exercise-9-6-17}
  \lean{PetersenLib.exercise9_6_17}
  \uses{rem:pet-ch9-exercise-9-6-16,def:pet-ch9-codazzi-harmonic-tensor}
  Let $E=TM$ in Exercise 9.6.16, so $\Omega^1(M,TM)=\operatorname{Hom}(TM,TM)$ consists of
  all $(1,1)$-tensors.
  \begin{enumerate}
    \item[(1)] Show $d^\nabla s=0$ if and only if $s$ is a Codazzi tensor.
    \item[(2)] The chapter suggests that whenever $s\in\Omega^p(M,E)$ satisfies
      $d^\nabla s=0$, there is a Bochner-type formula for $s$, and when $s$ is
      also divergence free with some curvature nonnegative, $s$ should be
      parallel. Develop such a theory in this generality.
    \item[(3)] Show that for a vector field $X$, $\nabla X$ is Codazzi if and
      only if $R(\cdot,\cdot)X=0$. Give an example where $\nabla X$ is Codazzi
      but $X$ is not parallel. Can a Bochner-type formula be established for
      exact tensors $\nabla X=d^\nabla X$ even when not closed?
  \end{enumerate}
  \source{petersen:page-0377}
\end{remark}

\begin{remark}[Exercise 9.6.18 (Thomas)]
  \label{rem:pet-ch9-exercise-9-6-18}
  \lean{PetersenLib.exercise9_6_18}
  \uses{rem:pet-ch9-exercise-9-6-17,def:pet-ch9-codazzi-harmonic-tensor}
  Show that in dimension $n\ge3$ the Gauss equations $\mathfrak{R}=S\wedge S$
  imply the Codazzi equations $d^\nabla S=0$, provided $\det S\ne0$. (Hint: use
  the second Bianchi identity, and study the linear map
  $\operatorname{Hom}(\Lambda^2V,V)\to\operatorname{Hom}(\Lambda^3V,\Lambda^2V)$, $T\mapsto T\wedge S$,
  for a linear map $S:V\to V$; this map is injective only when
  $\operatorname{rank}S\ge4$.)
  \source{petersen:page-0377}
\end{remark}

